\chapter{Discussion and Limitations}

\section{Main Findings}

\subsection{Lightweight Scene Understanding}

Among the evaluated deployment configurations, QAT SegFormer-B0 achieved the highest test mIoU, while QAT DeepLabV3+ with MobileNetV2 required the least GPU energy per image (Section~\ref{sec:quantization-results}). These configurations therefore represent the best observed accuracy-energy trade-offs: SegFormer-B0 favored segmentation quality, whereas DeepLabV3+ required less energy at the cost of lower accuracy. In the FP32 comparison, SegFormer-B0 exceeded DeepLabV3+ with MobileNetV2 by 11.48 percentage points of mIoU while requiring approximately 28\% more measured GPU energy per image. Against FP32 DeepLabV3 with ResNet34, the same FP32 SegFormer-B0 configuration achieved a relative mIoU advantage of approximately 2.4\% while consuming approximately 79\% less energy. These comparisons show that a compact Transformer configuration can provide competitive efficiency alongside accurate segmentation, although its advantage depends on the convolutional model considered. The preferred configuration consequently depends on the segmentation quality required by the application and the available energy budget.

The selected QAT configurations reduced measured GPU energy per image relative to their corresponding FP32 baselines by approximately 65\% for SegFormer-B0 and 82\% for DeepLabV3+ with MobileNetV2. Within this QAT comparison, SegFormer-B0 retained an mIoU advantage of 11.30 percentage points over DeepLabV3+ with MobileNetV2, but required approximately 153\% more energy. Its INT8 conversion was restricted to convolutions, leaving many parameters and the attention computations in FP16 to limit the risk of numerical degradation. This selective scope may partly explain the remaining inference cost. QAT also increased mIoU relative to FP32 by approximately 1.8\% and 2.4\% for SegFormer-B0 and DeepLabV3+, respectively, although these improvements resulted from a procedure that included fine-tuning and distillation. PTQ instead essentially preserved SegFormer-B0's accuracy but reduced DeepLabV3+ mIoU by approximately 7.7\% relative to FP32, confirming that accuracy preservation must be verified for each configuration. Whether the remaining energy difference is acceptable also depends on inference frequency and the available energy.

The remaining inference cost also illustrates why parameter count alone is not a reliable efficiency indicator. SegFormer-B0's attention operations and intermediate feature maps introduce computation and storage requirements that its lower parameter count does not capture, consistent with its higher measured memory use and energy consumption relative to DeepLabV3+.
The T4 benchmarks establish useful comparisons under the tested conditions, but deployment decisions require measurements on the intended hardware. The reported energy estimates cover the GPU, including idle consumption, and exclude CPU processing and preprocessing, which must also be considered when assessing a complete perception pipeline.

Beyond these resource considerations, the separate MarsSeg experiments provide an additional reason to consider SegFormer-B0. The evaluated configuration transferred better between MSL and MER than the convolutional alternatives, despite U-Net achieving the highest in-domain accuracy. These findings concern separately trained MarsSeg models rather than the quantized S5Mars deployments, but they indicate that model selection should also account for changes in acquisition conditions.

The training ablations likewise showed that additional variation and class balancing did not necessarily improve segmentation. The tested geometric augmentation reduced accuracy relative to the unaugmented reference, possibly by altering useful orientation and scene-layout cues. Oversampling also underperformed standard sampling, while uniform Dice weights partly recovered accuracy within the oversampled runs. This pattern may reflect overlapping emphasis on rare classes from sampling and weighted loss, or reduced observation diversity when sampling with replacement. These findings motivated keeping the tested augmentation policy and oversampling disabled in the reference configuration, while suggesting that alternative augmentation strategies should be tailored to semantic segmentation and better reflect the acquisition conditions and visual characteristics of Martian scenes.

\subsection{Synthetic Damage Inspection}
\label{sec:discussion-damage-inspection}

Within the synthetic inspection setting, the results support different choices according to whether detection quality or inference cost is given priority (Section~\ref{sec:damage-model-comparison}). PC-L provided the strongest image discrimination and the highest Pixel AP, but required more resources than SSN, TG-384, and TG-512. SSN retained stronger image-level detection than TG-384 and TG-512 at an intermediate cost, whereas TG-384 offered the lowest latency, memory use, and energy consumption. Within the latter family, TG-384 favored efficiency and TG-512 favored localization, with little change in image-level discrimination.

These trade-offs also expose a distinction between identifying a damaged observation and delineating its defect. TG-384 and TG-512 localized holes better than SSN according to Pixel AP and AUPRO, despite their weaker image discrimination. A high image score may therefore provide useful evidence of damage without a correspondingly precise map. Changing the SSN score formulation improved discrimination without modifying the map itself, showing that aggregation of local responses is part of the detector design. Detection and localization must consequently be assessed separately, with score selection matched to the spatial extent of the target anomaly.

Both objectives also depend on the acquisition setting. The study uses a single Curiosity asset and targets visible wheel perforations, with the main comparison restricted to pose A. This controlled viewpoint was chosen to provide a more consistent reference for normal appearance, while lighting and surface wear still vary. For SSN, the cropped pose A setting achieved higher Image AUROC and Pixel AP than the alternatives using all poses on validation, and this advantage was also observed on test (Section~\ref{sec:damage-supersimplenet-ablation}). These findings support the selected setting and cropping also concentrates the available resolution on the wheel, consistent with the observed localization gains.

Even with inputs focused on the wheel, localization still shows limited pixel precision and false positive responses on normal regions. High Pixel AUROC or AUPRO should therefore not be interpreted as evidence that the anomaly maps delineate holes accurately. The emphasis on compact models and moderate input resolutions may contribute by limiting the spatial detail used for inspection. Image discrimination provides more encouraging evidence for ranking damaged observations above normal ones, but the aggregate metrics do not establish an operational decision threshold or an acceptable false positive rate.

\section{Sim-to-Real Gap}

Transfer to real imagery introduced an additional difficulty. A preliminary qualitative check on a few wheel images obtained from NASA's website did not reproduce the quality observed on the synthetic data. This is consistent with a mismatch between the normal appearance learned from the renders and that of real inspection images, because differences in viewpoint, illumination, and surface texture may produce anomalous responses even in regions without visible holes. The few examined examples do not support a quantitative assessment of sim-to-real performance, but they suggest that learning this synthetic normality alone is insufficient for reliable real image inspection. Addressing this gap would require representative real wheel data or substantially closer matching of the synthetic scenes to real acquisition conditions, followed by independent validation on real images.
